% template.tex, dated April 5 2013
% This is a template file for Annual Reviews 1 column Journals
%
% Compilation using ar-1col.cls' - version 1.0, Aptara Inc.
% (c) 2013 AR
%
% Steps to compile: latex latex latex
%
% For tracking purposes => this is v1.0 - Apr. 2013

\documentclass{ar-1col}
\usepackage{url}
\usepackage{gb4e}
\usepackage{natbib}

\setcounter{secnumdepth}{4}

% Metadata Information
\jname{Xxxx. Xxx. Xxx. Xxx.}
\jvol{AA}
\jyear{YYYY}
\doi{10.1146/((please add article doi))}


% Document starts
\begin{document}

% Page header
\markboth{Author et al.}{Short title}

% Title
\title{Title: Subtitle}


%Authors, affiliations address.
\author{Author B. Authorone,$^1$ Firstname C. Authortwo,$^2$ and D. Name Authorthree$^3$
\affil{$^1$Department/Institute, University, City, Country, Postal code; email: author@email.edu}
\affil{$^2$Department/Institute, University, City, Country, Postal code}
\affil{$^3$Department/Institute, University, City, Country, Postal code}}

%Abstract
\begin{abstract}
Abstract text, approximately 150 words. 
\end{abstract}

%Keywords, etc.
\begin{keywords}
keywords, separated by comma, no full stop, lowercase
\end{keywords}
\maketitle

%Table of Contents
\tableofcontents


% Heading 1
\section{INTRODUCTION}
Please begin the main text of your article here. 



%Heading 1
\section{FIRST-LEVEL HEADING}
This is dummy text. 
% Heading 2
\subsection{Second-Level Heading}
This is dummy text. This is dummy text. This is dummy text. This is dummy text.

% Heading 3
\subsubsection{Third-Level Heading}
This is dummy text. This is dummy text. This is dummy text. This is dummy text. 

% Heading 4
\paragraph{Fourth-Level Heading} Fourth-level headings are placed as part of the paragraph.

%Example of a Figure
\section{ELEMENTS\ OF\ THE\ MANUSCRIPT}
All numbered items mentioned in text must be labeled as one of the following:

\begin{enumerate}
\item Figures (all images, including photos, drawings, maps, and trees)
\item Tables
\item Trees
\item Examples* (for numbered sentences, logical forms, etc.)
\end{enumerate}
*For numbered items that are not figures, tables, or trees, we suggest using the label ?example? (or ?logical form,? ?utterance,? etc.). 
 
\subsection{Figures}Figures should be cited in the main text in chronological order. This is dummy text with a citation to the first figure (\textbf{Figure \ref{fig1}}). Citations to \textbf{Figure \ref{fig1}} (and other figures) will be bold. 


\begin{figure}[h]
\includegraphics[width=3in]{SampleFigure}
\caption{Figure caption with descriptions of parts a and b}
\label{fig1}
\end{figure}

% Example of a Table
\subsection{Tables} Tables should also be cited in the main text in chronological order (\textbf {Table \ref{tab1}}).

\begin{table}[h]
\tabcolsep7.5pt
\caption{Table caption}
\label{tab1}
\begin{center}
\begin{tabular}{@{}l|c|c|c|c@{}}
\hline
Head 1 &&&&Head 5\\
{(}units)$^{a}$ &Head 2 &Head 3 &Head 4 &{(}units)\\
\hline
Column 1 &Column 2 &Column3$^{b}$ &Column4 &Column\\
Column 1 &Column 2 &Column3 &Column4 &Column\\
Column 1 &Column 2 &Column3 &Column4 &Column\\
Column 1 &Column 2 &Column3 &Column4 &Column\\
\hline
\end{tabular}
\end{center}
\begin{tabnote}
$^{a}$Table footnote; $^{b}$second table footnote.
\end{tabnote}
\end{table}

% Example of lists
\subsection{Lists} Here is an example of a numbered list:
\begin{enumerate}
\item List entry number 1,
\item List entry number 2,
\item List entry number 3,\item List entry number 4, and
\item List entry number 5.
\end{enumerate}

\subsection{Extracts}
\noindent Here is an example of a extract:
\begin{extract}
This is an example of a quote or extract.
This is an example of a quote or extract. This is an example of a quote or extract.
This is an example of a quote or extract. This is an example of a quote or extract.
This is an example of a quote or extract. 



\end{extract}

\subsection{Sidebars and Margin Notes}
See sidebar example above, and margin note example to the left.
% Margin Note
\begin{marginnote}[]
\entry{Term A}{definition}
\entry{Term B}{definition}
\entry{Term C}{defintion}
\end{marginnote}

\begin{textbox}[]\section{SIDEBARS}
Sidebar text goes here.
\subsection{Sidebar Second-Level Heading}
More text goes here.\subsubsection{Sidebar third-level heading}
Text goes here.\end{textbox}



\subsection{Equations}
% Example of a single-line equation
\begin{equation}
a = b \ {(Single\ Equation\ Numbered)}
\end{equation}
%Example of multiple-line equation
Equations can also be multiple lines as shown in Equations 2 and 3.
\begin{eqnarray}
c = 0 \ {(Multiple\  Lines, \ Numbered)}\\
ac = 0 \ {(Multiple \ Lines, \ Numbered)}
\end{eqnarray}

\subsection{Linguistic Examples}
\begin{exe}
\exi{(1a)}\label{Johna}      
John gave Mary a flower.
\exi{(1b)}\label{Johnb} 
John gave a flower to Mary.
\end{exe}

or

\begin{exe}
\exi{(3a)}\label{bear}     
\gll Pekka pel\"astyi karhusta.\\
Pekka {became afraid} bear.ELA\\
\trans `Pekka became afraid because of the/a bear.'
\end{exe}

% Summary Points
\begin{summary}[SUMMARY POINTS]
\begin{enumerate}
\item Summary point 1. These should be full sentences.
\item Summary point 2. These should be full sentences.
\item Summary point 3. These should be full sentences.
\item Summary point 4. These should be full sentences.
\end{enumerate}
\end{summary}

% Future Issues
\begin{issues}[FUTURE ISSUES]
\begin{enumerate}
\item Future issue 1. These should be full sentences.
\item Future issue 2. These should be full sentences.
\item Future issue 3. These should be full sentences.
\item Future issue 4. These should be full sentences.
\end{enumerate}
\end{issues}

%Disclosure
\section*{DISCLOSURE STATEMENT}
If the authors have noting to disclose, the following statement will be used: The authors are not aware of any affiliations, memberships, funding, or financial holdings that
might be perceived as affecting the objectivity of this review. 

% Acknowledgements
\section*{ACKNOWLEDGMENTS}
Acknowledgements, general annotations, funding.

% References
%
% Margin notes within bibliography
\section*{LITERATURE\ CITED}

To download the appropriate bibliography style file, please see \url{http://www.annualreviews.org/page/authors/author-instructions/preparing/latex}. 



\noindent
Please see the Style Guide document for instructions on preparing your Literature Cited.

The citations should be listed in alphabetical order, with titles. For example:





\begin{thebibliography}{00}

\bibitem[{Abrus\'an \& Szendr\H{o}i (2013)}]{AbrusanSzendroi2013}
Abrus\'an M, Szendr\H{o}i K. 2013. Experimenting with the king of
{F}rance: topics, verifiability, and definite descriptions.
\textit{Semant. Pragmat.} 6:1--43

\bibitem[{Abusch (2002)}]{Abusch2002}
Abusch D. 2002. Lexical alternatives as a source of pragmatic presuppositions. In \textit{Proceedings of the 12th Semantics and Linguistic Theory Conference}
({\textit{SALT}} \textit{12}), ed. B~Jackson. Ithaca, NY: CLC\ Publ.

\bibitem[{Abusch (2010)}]{Abusch2010}
Abusch D. 2010. Presupposition triggering from alternatives.
\textit{J. Semant.} 27:37--80

\bibitem[{Amaral \& Cummins (2015)}]{AmaralCummins2015}
Amaral P, Cummins C. 2015. A cross-linguistic study on information backgrounding and presupposition projection. In \textit{Experimental Perspectives on
Presuppositions}, ed. F~Schwarz, pp. 157--72. Berlin:\ Springer


\bibitem[{Beaver (2001)}]{Beaver2001}
Beaver D. 2001. {\it Presupposition and Assertion in Dynamic
Semantics}. Stanford, CA: Cent. Study Lang. Inf.

\bibitem[{Beaver \& Geurts (2012)}]{BeaverGeurts2012}
Beaver D, Geurts B. 2012. Presupposition. In \textit{Semantics: An International
Handbook of Natural Language Meaning}, Vol. 3, ed. C~Maienborn,
K~{von Heusinger}, P~Portner, pp. 2432--60. Berlin: de
Gruyter

\bibitem[{Beaver \& Krahmer (2001)}]{BeaverKrahmer2001}
Beaver D, Krahmer E. 2001. Presupposition and partiality: back to
the future. \textit{J. Logic Lang. Inform.} 10:147--82

\bibitem[{Beaver \& Zeevat (2007)}]{BeaverZeevat2007}
Beaver D, Zeevat H. 2007. Accommodation. In \textit{Oxford Handbook of Linguistic
Interfaces}, ed. G~Ramchand, C~Reiss, pp. 502--38. Oxford, UK:\ Oxford
Univ. Press

\bibitem[{Bill et~al. (2014)}]{Billetal2014}
Bill C, Romoli J, Schwarz F, Crain S. 2014. \textit{Indirect scalar
implicatures are neither scalar implicatures nor presuppositions}
(\textit{or both}). Presented at Ann. CUNY\ Conf. Hum. Sentence Process., 27th, Columbus, Ohio

\bibitem[{Bott \& Noveck (2004)}]{BottNoveck2004}
Bott L, Noveck IA. 2004. Some utterances are underinformative: the
onset and time course of scalar inferences. \textit{J. Mem. Lang.}
51:437--57

\bibitem[{Breheny, Ferguson \& Katsos (2013)}]{Brehenyetal2013}
Breheny R, Ferguson HJ, Katsos N. 2013. Investigating the
timecourse of accessing conversational implicatures during
incremental sentence interpretation. \textit{Lang. Cogn. Process.}
28:443--67




\end{thebibliography}


\end{document}
